\documentclass[11pt]{article}

\usepackage[utf8]{inputenc}
\usepackage[T1]{fontenc}
\usepackage{amsmath,amssymb}
\usepackage{graphicx}
\usepackage{booktabs}
\usepackage{hyperref}
\usepackage[margin=1in]{geometry}
\usepackage{natbib}
\usepackage{longtable}

% Graceful handling of missing figure files: if file does not exist, show placeholder text
\newcommand{\includefigure}[2][width=\textwidth]{%
  \IfFileExists{#2}{%
    \includegraphics[#1]{#2}%
  }{%
    \fbox{\parbox{0.9\textwidth}{\centering\vspace{2em}%
    \textit{Figure not available.}\par\vspace{0.5em}%
    {\small Regenerate ancestry comparison plots by running
    \texttt{scripts/validation/ancestry/plot\_ancestry\_comparison.py}}\par%
    \vspace{2em}}}%
  }%
}

\title{HPRCv2-IBD: Supplementary Materials}
\author{Franco et~al.}
\date{\today}

\begin{document}

\maketitle

\appendix

\section{Parameter Sensitivity Analysis}
\label{sec:supp_sensitivity}

\subsection{Window Size}

The window size $W$ controls the trade-off between spatial resolution and statistical power for IBD detection. Smaller windows increase boundary precision but reduce the per-window signal-to-noise ratio (higher variance in identity estimates). The default of 10~kb was selected as a balance point: it provides approximately 8--9 informative SNPs per window in European populations (at $\pi = 0.00085$), sufficient for meaningful per-window identity estimates while maintaining $\sim$13,000 windows per chromosome for robust HMM inference. All validation results in this paper use the 10~kb default; systematic evaluation across window sizes (1--20~kb) is deferred to future work.

\subsection{Transition Parameters}

\subsubsection{Expected Segment Length and Entry Probability}

The expected IBD segment length $L$ (in windows) parameterizes the ``stickiness'' of the IBD state, while the entry probability $p_{\text{enter}}$ controls overall sensitivity. However, as demonstrated in the main text (Performance Ceiling Analysis), the IBD detection bottleneck lies entirely in the emission model rather than in transitions. On chr20 validation data, three different transition models (standard, distance-aware, genetic-map-aware) produce identical segment calls (Table~\ref{tab:distance_aware}). Baum-Welch training, which re-estimates both emission and transition parameters from data, provides the dominant improvement ($+$110\% F1). Because Baum-Welch learns effective transition parameters from the data, the initial settings of $L$ and $p_{\text{enter}}$ are less critical than the emission model quality. We use defaults of $L = 50$ windows and $p_{\text{enter}} = 10^{-4}$.

\subsection{Baum-Welch Iterations}

Baum-Welch training is the most impactful modeling choice for IBD detection, providing a $+$110\% improvement in F1 compared to no training (BW=0) at matched parameters (identity floor = 0.9, min-len = 2~Mb).

\begin{table}[ht]
\centering
\caption{Effect of Baum-Welch iterations on IBD concordance (chr20, floor = 0.9, min-len = 2 Mb).}
\label{tab:bw_iterations}
\begin{tabular}{lcccc}
\hline
\textbf{BW iterations} & \textbf{Segments} & \textbf{F1} & \textbf{$\Delta$ vs.\ BW=0} \\
\hline
0 & 7 & 13.4\% & --- \\
5 & 12 & 26.1\% & $+$95\% \\
\textbf{10} & \textbf{12} & \textbf{28.3\%} & \textbf{$+$110\%} \\
20 & 12 & 28.3\% & $+$110\% \\
\hline
\end{tabular}
\end{table}

On this dataset, convergence is reached by 10 iterations, with no further improvement at 20. However, convergence rate analysis (see Parameter Estimation in the main text) shows that the asymptotic rate $r \approx 0.9$ means 10 iterations leave $\sim$35\% residual parameter error in the general case; we therefore recommend a default of 20 iterations for robustness across populations and datasets. Note that for ancestry inference, Baum-Welch training is \emph{counterproductive}, degrading concordance by 2--3 percentage points (Table~\ref{tab:ancestry_bw_iters}), because the weak pangenome similarity signal does not support reliable data-driven transition estimation.

\subsection{Distance-Aware vs.\ Standard Inference}

We compared three transition models on the chr20 validation dataset (44 sample pairs, identity floor = 0.9, min-len = 2 Mb, Baum-Welch = 10).

\begin{table}[ht]
\centering
\caption{Effect of transition model on IBD concordance with hap-ibd (chr20, floor = 0.9, 2 Mb).}
\label{tab:distance_aware}
\begin{tabular}{lccccc}
\hline
\textbf{Transition model} & \textbf{Segments} & \textbf{F1} & \textbf{Jaccard} & \textbf{Precision} & \textbf{Recall} \\
\hline
Standard (uniform) & 12 & 28.3\% & 22.9\% & 31.7\% & 27.0\% \\
Distance-aware (physical) & 12 & 28.3\% & 22.9\% & 31.7\% & 27.0\% \\
Genetic-map-aware (Haldane) & 12 & 28.3\% & 22.9\% & 31.7\% & 27.0\% \\
\hline
\end{tabular}
\end{table}

All three transition models produce identical segment calls and identical concordance metrics. This demonstrates that the performance bottleneck for pangenome-based IBD detection lies entirely in the emission model: when the emission likelihoods for IBD and non-IBD states are nearly equal (identity 0.997 vs.\ 0.9997), transition adjustments---whether by physical distance, genetic distance, or recombination rate---have no practical effect on the Viterbi path or forward-backward posteriors.

\subsection{LOD Score Threshold}

Segments can be filtered by minimum LOD score. In validation, posterior threshold had no effect on the best configuration because all 12 detected segments at the optimal settings (floor=0.9, 2~Mb, BW=10) already have $P(\text{IBD}) > 0.8$ and high LOD scores ($>$100). For the hap-ibd reference set, a \texttt{validate-min-lod} threshold of 3.0 was found optimal: lower values include weak hap-ibd segments that represent borderline detections even in VCF space, inflating the false negative count without improving validation quality.

\subsection{Posterior Threshold}

The posterior threshold filters segments by minimum mean $P(\text{IBD})$. At the optimal IBD configuration (floor=0.9, 2~Mb, BW=10), all detected segments have mean posterior $> 0.8$, so thresholds up to 0.8 have no effect on the result set. This occurs because Baum-Welch training sharpens the emission distributions sufficiently that surviving segments have unambiguous posterior support. The posterior threshold is most useful at relaxed settings (e.g., 1~Mb minimum length) where it can remove a small fraction of marginal detections.

\section{Ancestry Validation: Parameter Sensitivity}
\label{sec:supp_ancestry_sensitivity}

All ancestry parameter sensitivity analyses were performed on HG00733 (Puerto Rican), chromosome 20, using a matched reference panel of 30 EUR + 49 AFR individuals. Concordance is measured against RFMix v2 as described in the main text (Ancestry Validation). Unless otherwise noted, default parameters are: emission model = mean, $\tau = 0.03$, $p_{\text{switch}} = 0.01$, BW iterations = 0, smoothing = 3.

\subsection{Temperature}

\begin{table}[ht]
\centering
\caption{Effect of softmax temperature on ancestry concordance with RFMix (HG00733, chr20).}
\label{tab:temperature_sensitivity}
\begin{tabular}{lcccc}
\hline
\textbf{Temperature} $\boldsymbol{\tau}$ & \textbf{Concordance} & \textbf{EUR F1} & \textbf{AFR F1} & \textbf{Runtime (s)} \\
\hline
0.01 & 0.887 & 0.940 & 0.074 & 4.9 \\
0.03 (default) & 0.887 & 0.940 & 0.033 & 4.9 \\
0.10 & 0.890 & 0.941 & 0.030 & 4.9 \\
0.50 & 0.876 & 0.934 & 0.026 & 5.2 \\
1.0 & 0.872 & 0.931 & 0.027 & 4.9 \\
5.0 & 0.870 & 0.930 & 0.027 & 5.1 \\
\hline
\end{tabular}
\end{table}

Temperature has minimal effect on concordance (range: 0.870--0.890). Lower temperatures produce slightly sharper posteriors and marginally better concordance, but the effect is small compared to emission model choice.

\subsection{Switch Probability}

\begin{table}[ht]
\centering
\caption{Effect of switch probability on ancestry concordance with RFMix (HG00733, chr20).}
\label{tab:switch_sensitivity}
\begin{tabular}{lccc}
\hline
$\boldsymbol{p_{\text{switch}}}$ & \textbf{Concordance} & \textbf{EUR F1} & \textbf{AFR F1} \\
\hline
0.001 & 0.909 & 0.952 & 0.000 \\
0.005 & 0.897 & 0.945 & 0.018 \\
0.01 & 0.887 & 0.940 & 0.033 \\
0.05 & 0.856 & 0.922 & 0.085 \\
\hline
\end{tabular}
\end{table}

Lower switch probabilities improve raw concordance but prevent detection of minority ancestry tracts. At $p_{\text{switch}} = 0.001$, AFR F1 drops to zero because the model effectively never switches away from the majority EUR state. The trade-off between concordance and minority detection is fundamental: tuning for one degrades the other.

\subsection{Emission Model Comparison}

\begin{table}[ht]
\centering
\caption{Ancestry concordance by emission aggregation model (HG00733, chr20).}
\label{tab:emission_model_cv}
\begin{tabular}{lcccc}
\hline
\textbf{Model} & \textbf{Concordance} & \textbf{EUR F1} & \textbf{AFR F1} & \textbf{$\Delta$ vs.\ mean} \\
\hline
Max & 0.932 & 0.965 & 0.158 & $+$4.5 pp \\
Mean & 0.887 & 0.940 & 0.033 & --- \\
Median & 0.864 & 0.926 & 0.086 & $-$2.3 pp \\
\hline
\end{tabular}
\end{table}

The \texttt{max} model is the single most impactful parameter, improving concordance by 4.5 percentage points over \texttt{mean}. It also provides the best minority ancestry detection (AFR F1 = 0.158 vs.\ 0.033 for mean). The max model uses the highest similarity to any reference haplotype within each population, effectively performing a nearest-neighbor ancestry assignment.

\subsection{Baum-Welch Iterations}

\begin{table}[ht]
\centering
\caption{Effect of Baum-Welch iterations on ancestry concordance (HG00733, chr20, mean model).}
\label{tab:ancestry_bw_iters}
\begin{tabular}{lcccc}
\hline
\textbf{Iterations} & \textbf{Trained $p_{\text{switch}}$} & \textbf{Concordance} & \textbf{EUR F1} & \textbf{AFR F1} \\
\hline
0 (disabled) & --- & 0.887 & 0.940 & 0.033 \\
5 & 0.032 & 0.863 & 0.926 & 0.075 \\
10 & 0.046 & 0.858 & 0.923 & 0.083 \\
\hline
\end{tabular}
\end{table}

Baum-Welch training \emph{degrades} ancestry concordance by 2--3 percentage points. The EM algorithm overestimates the switch probability (from 0.01 to 0.03--0.05), producing noisier ancestry calls. This is because the pangenome similarity signal is not strong enough to support data-driven transition estimation for local ancestry; prior-driven parameters informed by the expected admixture complexity are more robust.

\subsection{Smoothing}

\begin{table}[ht]
\centering
\caption{Effect of smoothing on ancestry concordance (HG00733, chr20, mean model).}
\label{tab:smoothing_sensitivity}
\begin{tabular}{lccc}
\hline
\textbf{Min windows} & \textbf{Concordance} & \textbf{EUR F1} & \textbf{AFR F1} \\
\hline
0 & 0.887 & 0.940 & 0.033 \\
3 & 0.887 & 0.940 & 0.033 \\
5 & 0.887 & 0.940 & 0.033 \\
10 & 0.887 & 0.940 & 0.033 \\
\hline
\end{tabular}
\end{table}

Smoothing has no measurable effect, indicating the HMM already produces smooth state assignments on this data without post-hoc filtering.

\subsection{Auto-Estimated Parameters}

\begin{table}[ht]
\centering
\caption{Ancestry concordance with auto-estimated parameters (HG00733, chr20).}
\label{tab:ancestry_auto}
\begin{tabular}{llcccc}
\hline
\textbf{Emission} & \textbf{BW} & \textbf{Est.\ $\tau$} & \textbf{Est.\ $p_{\text{sw}}$} & \textbf{Concordance} & \textbf{AFR F1} \\
\hline
max & 0 & 0.01 & 0.0019 & 0.929 & 0.015 \\
max & 10 & 0.01 & 0.0019 & 0.904 & 0.166 \\
mean & 0 & 0.01 & 0.0019 & 0.899 & 0.032 \\
mean & 10 & 0.01 & 0.0019 & 0.838 & 0.096 \\
median & 0 & 0.01 & 0.0019 & 0.880 & 0.015 \\
median & 10 & 0.01 & 0.0019 & 0.806 & 0.099 \\
\hline
\end{tabular}
\end{table}

Auto-estimated parameters achieve near-optimal results: the best auto configuration (max, no BW) reaches 92.9\% concordance, only 0.3 pp below the manually optimized 93.2\%. This validates the parameter estimation procedure for practical use.

\subsection{Summary of Ancestry Sensitivity}
\label{sec:ancestry_sensitivity_summary}

\begin{table}[ht]
\centering
\caption{Recommended ancestry parameters and their effect size.}
\label{tab:ancestry_recommendations}
\begin{tabular}{llcc}
\hline
\textbf{Parameter} & \textbf{Recommended} & \textbf{Concordance range} & \textbf{Most impactful?} \\
\hline
Emission model & max & 0.864--0.932 & \textbf{Yes} (6.8 pp) \\
Temperature & 0.03 & 0.870--0.890 & No (2.0 pp) \\
Switch prob. & 0.001--0.01 & 0.856--0.909 & Moderate (5.3 pp) \\
BW iterations & 0 & 0.858--0.887 & Negative effect \\
Smoothing & 0--3 & No change & None \\
\hline
\end{tabular}
\end{table}

\subsection{Identity Floor Parameter}
\label{sec:supp_identity_floor}

The identity floor parameter removes windows below a threshold before HMM inference, addressing the bimodal pangenome identity distribution (see Pangenome Identity in the main text).

\begin{table}[ht]
\centering
\caption{Effect of identity floor on IBD concordance with hap-ibd (chr20, Baum-Welch = 10).}
\label{tab:identity_floor}
\begin{tabular}{llccccc}
\hline
\textbf{Floor} & \textbf{Min-len} & \textbf{Segments} & \textbf{F1} & \textbf{Jaccard} & \textbf{Precision} & \textbf{Recall} \\
\hline
\multicolumn{7}{l}{\textit{With 2 Mb minimum length (optimal):}} \\
0.0 (none) & 2 Mb & 12 & 25.3\% & 19.6\% & --- & --- \\
0.5 & 2 Mb & 12 & 25.7\% & 20.1\% & --- & --- \\
\textbf{0.9} & \textbf{2 Mb} & \textbf{12} & \textbf{28.3\%} & \textbf{22.9\%} & \textbf{31.7\%} & \textbf{27.0\%} \\
0.95 & 2 Mb & 12 & 28.3\% & 22.9\% & 31.7\% & 27.0\% \\
0.98 & 2 Mb & 11 & 25.3\% & 19.7\% & --- & --- \\
0.99 & 2 Mb & 10 & 23.9\% & 19.2\% & --- & --- \\
\hline
\multicolumn{7}{l}{\textit{With 1 Mb minimum length:}} \\
0.0 (none) & 1 Mb & 220 & 15.2\% & 11.6\% & 14.6\% & 19.3\% \\
0.98 & 1 Mb & 140 & 18.3\% & 14.7\% & 18.1\% & 21.5\% \\
0.995 & 1 Mb & 38 & 16.7\% & 14.0\% & 20.1\% & 15.1\% \\
\hline
\end{tabular}
\end{table}

The optimal configuration combines an identity floor of 0.9 with a 2 Mb minimum segment length, improving F1 from 10.7\% (default) to 28.3\%. The floor of 0.9--0.95 is optimal: lower values allow alignment-gap noise (many windows have identity near 0), while higher values (0.98--0.99) remove informative windows with identity between 0.9 and 0.98 that contribute to the HMM's ability to distinguish states. Baum-Welch training provides a $+$110\% F1 improvement at matched parameters (BW=10 vs.\ BW=0 at floor=0.9, 2 Mb).

\subsection{Recombination-Aware vs.\ Standard vs.\ Distance-Aware}

See Table~\ref{tab:distance_aware} for the comparison of all three transition models. On the chr20 validation data, all three models produce identical IBD calls. The transition model may become more impactful when: (a) the emission model is improved to provide stronger state discrimination, or (b) the analysis is extended to chromosomes with extreme recombination rate variation (e.g., chromosome 6 near the MHC locus). The current result demonstrates that transition refinement is a second-order effect relative to emission model quality.

\subsection{Coverage-Ratio Auxiliary Emission}
\label{sec:supp_coverage_ratio}

Alignment coverage ratio---defined as $\min(a_{\text{len}}, b_{\text{len}}) / \max(a_{\text{len}}, b_{\text{len}})$ for aligned segment lengths per haplotype pair---provides a genuinely independent signal from sequence identity (Pearson $r = 0.33$). IBD windows exhibit more symmetric alignment coverage ($\sim$0.996) than non-IBD windows ($\sim$0.859).

\begin{table}[ht]
\centering
\caption{Effect of coverage-ratio auxiliary emission on IBD concordance with hap-ibd (chr20, floor = 0.9, BW = 10).}
\label{tab:coverage_ratio}
\begin{tabular}{llccccc}
\hline
\textbf{Config} & \textbf{Min-len} & \textbf{Segments} & \textbf{F1} & \textbf{Jaccard} & \textbf{Concordance} \\
\hline
Baseline & 2 Mb & 12 & 8.74\% & 7.07\% & 94.0\% \\
Coverage & 2 Mb & 55 & 6.59\% & 4.77\% & 93.2\% \\
Baseline & 3 Mb & 7 & 6.63\% & 5.37\% & 93.5\% \\
\textbf{Coverage} & \textbf{3 Mb} & \textbf{12} & \textbf{8.64\%} & \textbf{6.58\%} & \textbf{94.1\%} \\
Baseline & 4 Mb & 3 & 2.99\% & 2.48\% & 93.3\% \\
Coverage & 4 Mb & 4 & 4.38\% & 3.66\% & 93.1\% \\
\hline
\end{tabular}
\end{table}

Coverage helps at strict minimum lengths (3--4 Mb) by detecting segments that identity alone misses, but at the standard 2 Mb threshold, it generates excessive false positives: 55 segments (vs.\ 12 baseline). Per-pair analysis reveals that the coverage model detects 3 additional true positive pairs but introduces 22 false-positive-only pairs, reducing pair-level precision from 0.60 to 0.29.

\subsection{Region Exclusion Filtering}
\label{sec:supp_region_exclusion}

Post-HMM region exclusion filtering blacklists known false positive hotspot regions (chr20:35--40 Mb and 45--50 Mb, accounting for 57\% of false positive segments).

\begin{table}[ht]
\centering
\caption{Effect of region exclusion on IBD concordance (chr20, floor = 0.9, 2 Mb, BW = 10).}
\label{tab:region_exclusion}
\begin{tabular}{lccccc}
\hline
\textbf{Config} & \textbf{Segments} & \textbf{F1} & \textbf{Jaccard} & \textbf{Precision} & \textbf{Recall} \\
\hline
Baseline (no exclusion) & 7 & 25.4\% & 21.6\% & 27.0\% & 24.9\% \\
Baseline + exclusion & 3 & 15.9\% & 14.1\% & 16.2\% & 15.9\% \\
Coverage (no exclusion) & 55 & 14.2\% & 10.3\% & 13.8\% & 17.6\% \\
Coverage + exclusion & 23 & 10.3\% & 7.4\% & 10.3\% & 12.5\% \\
\hline
\end{tabular}
\end{table}

Region exclusion consistently degrades results because the artifact regions contain real IBD segments (true positives) that are removed alongside false positives.

\subsection{Ancestry: Coverage-Ratio Auxiliary Feature}
\label{sec:supp_ancestry_coverage}

We tested alignment coverage ratio as an auxiliary emission feature for ancestry inference (HG00733, chr20).

\begin{table}[ht]
\centering
\caption{Effect of coverage-ratio auxiliary emission on ancestry concordance (HG00733, chr20, max emission).}
\label{tab:ancestry_coverage}
\begin{tabular}{lcccc}
\hline
\textbf{Config} & \textbf{Concordance} & \textbf{Hap1 AFR F1} & \textbf{Hap2 AFR F1} & \textbf{Mean AFR F1} \\
\hline
Baseline (no coverage) & 0.929 & 0.000 & 0.030 & 0.015 \\
Coverage $w$=0.5 & 0.896 & 0.037 & 0.003 & 0.020 \\
Coverage $w$=1.0 & 0.888 & 0.089 & 0.000 & 0.045 \\
Coverage $w$=2.0 & 0.874 & 0.068 & 0.008 & 0.038 \\
\hline
\end{tabular}
\end{table}

Coverage degrades overall concordance (3--6 pp) while producing inconsistent improvements across haplotypes. This further confirms that alignment-derived features cannot overcome the fundamental signal-to-noise limitation.

\subsection{Multi-Feature Oracle Ceiling Analysis}
\label{sec:supp_oracle}

To establish a definitive upper bound on the discriminative power of pangenome alignment features for ancestry, we tested all combinations of seven features (max/mean identity, max/mean Jaccard, max/mean coverage ratio, alignment count ratio) with oracle access to optimal classification thresholds using RFMix ground truth.

Every single-feature, two-feature, and three-feature oracle classifier achieved accuracy exactly equal to the majority-class baseline (95.0\% for haplotype 1, 93.4\% for haplotype 2) with zero AFR recall. Signal-to-noise ratios ranged from 0.012 (alignment count ratio) to 0.078 (max identity), all far below the threshold required for per-window discrimination.

This analysis establishes that the HMM's 93.2\% concordance is entirely attributable to temporal smoothing via the transition model, not to per-window signal detection. The per-window information content of all available pangenome alignment features is fundamentally insufficient for ancestry classification.

\subsection{Ancestry: Posterior Decoding vs.\ Viterbi}
\label{sec:supp_posterior_decoding}

Posterior decoding makes independent per-window ancestry assignments based on marginal posteriors, avoiding the path constraint of Viterbi decoding.

\begin{table}[ht]
\centering
\caption{Decoding method and normalization comparison for ancestry inference (HG00733, chr20, max emission, $\tau = 0.03$, $p_{\text{switch}} = 0.01$).}
\label{tab:decoding_comparison}
\begin{tabular}{llccc}
\hline
\textbf{Decoding} & \textbf{Normalization} & \textbf{Concordance} & \textbf{EUR F1} & \textbf{AFR F1} \\
\hline
Viterbi & Off & 0.932 & 0.965 & 0.158 \\
Posterior & Off & 0.932 & 0.965 & 0.158 \\
Viterbi & On & 0.103 & --- & 0.12 \\
Posterior & On & 0.103 & --- & 0.12 \\
\hline
\end{tabular}
\end{table}

Posterior decoding produces results identical to Viterbi (with and without normalization). This occurs because the forward-backward posteriors rarely favor the minority ancestry: even in true AFR regions, $P(\text{AFR})$ reaches a maximum of 0.67, with only 12 of 29 windows exceeding 0.5 in a representative 285 kb AFR tract. The signal is insufficient for the marginal posteriors to disagree with the MAP path.

\subsection{Ancestry: Emission Normalization}
\label{sec:supp_normalization}

Per-population emission normalization (z-score standardization before softmax) is catastrophic for ancestry concordance: 93.2\% $\rightarrow$ 10.3\%. The normalization removes the systematic EUR--AFR similarity difference that constitutes the discriminative signal. After z-score normalization, the populations have equal mean similarity by construction, and the model loses all ability to distinguish ancestries. This is a fundamental insight: unlike in other applications where normalization removes unwanted bias, here the ``bias'' (higher EUR similarity) \emph{is} the ancestry signal.

\subsection{Ancestry LOD Score Threshold}

Ancestry LOD scores quantify confidence in each segment's ancestry assignment as
\[
\text{LOD} = \sum_w \log_{10}\!\bigl(P(o_w \mid \text{assigned}) \;/\; P(o_w \mid \text{second-best})\bigr).
\]
Higher LOD thresholds remove ambiguous segments, which occur most commonly at ancestry transition boundaries and in populations with similar pangenome identity profiles. Given the low ancestry SNR (see Ancestry SNR Analysis in the main text), LOD filtering primarily affects short segments at ancestry boundaries, with limited impact on overall concordance because the majority of the chromosome is unambiguously assigned to the majority ancestry (EUR). LOD-based filtering is most valuable when computing admixture proportions, where removing ambiguous assignments near switch points reduces noise in the proportion estimates.

\subsection{Admixture Proportion Accuracy}

Admixture proportions are estimated from ancestry tract lengths (see Admixture Estimation in the main text). For HG00733 on chr20, our tool estimates approximately 94\% EUR and 6\% AFR ancestry, consistent with both the RFMix assignment and the expectation for an admixed Puerto Rican individual on this chromosome. Across all 16 AMR samples (see Multi-Sample Ancestry in the main text), our tool detects 0.4--7.0\% AFR ancestry; RFMix assigns 0.0--5.7\%. Our tool systematically over-detects AFR ancestry by $\sim$2 pp in samples that RFMix calls 100\% EUR, due to noise-driven HMM switches (Section~\ref{sec:supp_ancestry_segments}). Quantitative validation of admixture proportion accuracy across a wider range of true proportions (e.g., 50/50, 70/30, 90/10) would require simulated admixed genomes with known ground truth; this is deferred to future work.

\section{Scaled IBD Validation}
\label{sec:supp_scaled_ibd}

IBD validation was scaled from the initial 44 sample pairs to all 2,850 possible pairs among 76 HPRC samples (30 EUR + 30 AFR + 16 AMR) on chr20. Table~\ref{tab:scaled_summary} summarizes the population-stratified results.

\begin{table}[ht]
\centering
\caption{Scaled IBD validation summary: all 2,850 HPRC pairs on chr20 (floor = 0.9, 2 Mb, BW = 10).}
\label{tab:scaled_summary}
\begin{tabular}{lccccccc}
\hline
\textbf{Population} & \textbf{Pairs} & \textbf{Our segs} & \textbf{TP} & \textbf{FP} & \textbf{FN} & \textbf{Mean seg F1} & \textbf{Concordance} \\
\hline
EUR & 435 & 89 & 8 & 79 & 1 & 0.443 & 0.930 \\
AFR & 435 & 28 & 2 & 26 & 3 & 0.750 & 0.963 \\
AMR & 120 & 42 & 1 & 41 & 1 & 0.180 & 0.946 \\
AMR--EUR & 480 & 90 & 0 & 90 & 0 & --- & --- \\
AFR--EUR & 900 & 8 & 0 & 8 & 0 & --- & --- \\
AFR--AMR & 480 & 3 & 0 & 3 & 0 & --- & --- \\
\hline
\textbf{Total} & \textbf{2,850} & \textbf{293} & \textbf{11} & \textbf{247} & \textbf{5} & \textbf{0.475} & \textbf{0.938} \\
\hline
\end{tabular}
\end{table}

The full-scale analysis detected 293 IBD segments across 258 pairs, compared to 17 hap-ibd segments (LOD $\geq$ 3) across 16 pairs. Eleven pairs are true positives (detected by both methods), with mean segment F1 = 0.475 for TP pairs. False positive segments are concentrated in two genomic hotspots: chr20:35--40 Mb (51\% of FPs) and chr20:45--50 Mb (25\%). AMR--EUR cross-population pairs contribute 90/247 (36\%) of false positive pairs, reflecting shared European ancestry tracts in admixed samples. The consistency of false positive patterns between the 44-pair and 2,850-pair analyses confirms that these are properties of the pangenome alignment, not artifacts of sample selection.

Post-hoc threshold analysis shows that raising the minimum segment length to 3.0 Mb or requiring LOD $\geq$ 300 improves segment-level precision from 4.8\% to 35--46\% while retaining approximately half of true positive segments (see Threshold Analysis in the main text).

\subsection{Ancestry Segment-Level and Multi-Sample Analysis}
\label{sec:supp_ancestry_segments}

In addition to per-window concordance, we performed segment-level analysis comparing our ancestry calls with RFMix v2 on HG00733 chr20 (Table~\ref{tab:ancestry_segment_summary}).

\begin{table}[ht]
\centering
\caption{Ancestry segment-level comparison: our tool vs.\ RFMix v2 (HG00733, chr20).}
\label{tab:ancestry_segment_summary}
\small
\begin{tabular}{lccccc}
\hline
\textbf{Haplotype} & \textbf{Our segs} & \textbf{RFMix segs} & \textbf{Conc.} & \textbf{Bdy.\ offset (med.)} & \textbf{Discord.} \\
\hline
Hap1 & 26 & 63 & 94.6\% & 73 kb & 24 \\
Hap2 & 10 & 63 & 92.9\% & 83 kb & --- \\
\hline
\textbf{Combined} & \textbf{36} & \textbf{126} & \textbf{93.8\%} & --- & --- \\
\hline
\end{tabular}
\end{table}

Our tool produces substantially fewer segments (36 vs.\ 126 across both haplotypes) because it cannot resolve the short AFR ancestry tracts that RFMix detects. Of the 24 discordant regions on haplotype 1, 19 (79\%) are cases where our tool assigns EUR or AMR but RFMix assigns AFR, and 5 (21\%) are cases where our tool assigns AFR but RFMix assigns EUR. The discordant EUR/AMR-vs-AFR regions range from 30 kb to 424 kb (mean: 152 kb), confirming that short minority ancestry tracts below $\sim$500 kb are undetectable by our method.

Multi-sample ancestry validation across 15 AMR individuals (see Multi-Sample Ancestry in the main text) confirms that concordance with RFMix is stable (mean = 93.4\%, range 90.2--95.5\%). A systematic $\sim$2 pp over-detection of AFR ancestry occurs in samples where RFMix calls 100\% EUR, reflecting noise-driven HMM switches rather than true ancestry signal. Among samples with genuine admixture (RFMix AFR $>$ 3\%), AFR F1 ranges from 0.061 to 0.706, with higher values for individuals whose AFR tracts are long and concentrated rather than fragmented.

\section{Preliminary Cross-Chromosome Analysis: Chr15}
\label{sec:supp_chr15}

As a preliminary cross-chromosome analysis prior to the full chr10/11/12 validation (see IBD Ground Truth in the main text), we ran hap-ibd on the same 76 HPRC samples using a phased chr15 VCF (4,091 samples, 170,850 markers). The chr15 analysis motivated the multi-chromosome validation design by demonstrating that IBD hotspots are a recurring feature across chromosomes. Table~\ref{tab:chr15_overview} summarizes the chr15 hap-ibd ground truth.

\begin{table}[ht]
\centering
\caption{hap-ibd IBD segments on chr15 (76 HPRC samples, LOD $\geq$ 3), compared with chr20.}
\label{tab:chr15_overview}
\begin{tabular}{lccc}
\hline
\textbf{Metric} & \textbf{Chr20} & \textbf{Chr15} & \textbf{Chr15 (non-hotspot)} \\
\hline
Total segments (LOD $\geq$ 3) & 17 & 2,317 & 17 \\
Unique pairs & 16 & 1,476 & 16 \\
Mean segment length (Mb) & 3.60 & 2.75 & 1.80 \\
Mean LOD & 4.8 & 14.7 & 4.5 \\
\hline
\end{tabular}
\end{table}

A massive IBD hotspot at chr15:17.0--22.6 Mb (the chr15q11--q13 region~\cite{nicholls2001prader_willi}) accounts for 2,300 of 2,317 (99.3\%) LOD $\geq$ 3 segments. In this 5.6 Mb region, 1,472 of 2,850 possible sample pairs (51.6\%) share IBD---an order-of-magnitude higher IBD density than anywhere on chr20. The chr15q11--q13 region is known for complex genomic architecture including segmental duplications, imprinted genes, and frequent structural rearrangements, which may produce extended blocks of identity-by-state that inflate IBD detection.

Outside the hotspot, chr15 shows 17 IBD segments among 16 pairs---remarkably consistent with chr20's 17 segments among 16 pairs (Table~\ref{tab:chr15_overview}). Cross-chromosome pair overlap analysis shows that 11 of 16 chr20 IBD pairs also share IBD on chr15 (within the hotspot region), consistent with genome-wide IBD sharing among closely related individuals. Mean segment length outside the chr15 hotspot (1.80 Mb) is shorter than on chr20 (3.60 Mb), suggesting weaker or older shared ancestry on chr15 outside the complex region.

Preliminary pangenome IBD validation was performed on a 2.17 Mb region of chr15 (31.0--33.2 Mb), where 217 IBS windows were computed for all 2,850 sample pairs. Because the analyzed region is smaller than 2 Mb, a 1 Mb minimum segment length was used.

\begin{table}[ht]
\centering
\caption{Cross-chromosome IBD validation: chr15 region 2 (31.0--33.2 Mb) vs.\ chr20 (full chromosome).}
\label{tab:chr15_validation}
\begin{tabular}{lcc}
\hline
\textbf{Metric} & \textbf{Chr20} & \textbf{Chr15 (31--33 Mb)} \\
\hline
Our IBD segments & 293 & 39 \\
Our detected pairs & 258 & 39 \\
hap-ibd segments (LOD $\geq$ 3) & 17 & 7 \\
TP pairs & 11 & 1 \\
FP pairs & 247 & 38 \\
FN pairs & 5 & 6 \\
Pair precision & 4.3\% & 2.6\% \\
Pair recall & 68.8\% & 14.3\% \\
False positive rate (\% detected) & 95.7\% & 97.4\% \\
\hline
\end{tabular}
\end{table}

The false positive rate is remarkably consistent across chromosomes ($\sim$96--97\%), confirming that the high FP rate is a systematic measurement limitation rather than a chromosome-specific artifact. The single true positive pair (HG00272 vs.\ HG00323) has a strong hap-ibd signal (LOD = 12.2), while all 6 missed pairs have near-threshold LOD scores (3.0--4.1). This is consistent with the chr20 finding that only strong hap-ibd signals (LOD $>$ 5) are reliably detected by pangenome identity. The chr15 hotspot at chr15q11--q13 remains the most promising target for future validation: its mean LOD of 14.7 may yield better pangenome--hap-ibd concordance if the signal is the binding constraint, though the region's structural complexity could confound alignment quality.

\section{Extended Figures}
\label{sec:supp_figures}

\begin{enumerate}
    \item \textbf{Figure S1: Pangenome identity distribution.} Histogram of per-window sequence identity for a representative sample pair, showing the bimodal distribution (alignment gaps near 0, aligned regions near 1) and the effect of identity floor filtering.

    \item \textbf{Figure S2: Emission distribution fits.} Observed identity histogram (after floor filtering) with fitted 2-component Gaussian overlay for EUR, AFR, and EAS populations, illustrating the population-specific non-IBD/IBD separation.

    \item \textbf{Figure S3: Chromosome painting comparison.} Side-by-side ancestry painting of HG00733 chr20 from our tool and RFMix. Each row shows one haplotype; colors represent EUR (blue) and AFR (red) ancestry.

    \item \textbf{Figure S4: Ancestry parameter sensitivity heatmap.} Concordance as a function of emission model and switch probability, showing that \texttt{max} emission with low switch probability yields optimal concordance.

    \item \textbf{Figure S5: IBD concordance vs.\ segment length.} Scatter plot of per-pair Jaccard index vs.\ hap-ibd segment length, showing that concordance increases with segment length. Points colored by population pair type.

    \item \textbf{Figure S6: Ancestry posterior along chromosome.} Per-window posterior probabilities for EUR and AFR in HG00733 chr20, overlaid with RFMix ground truth, illustrating the weak posterior signal for minority ancestry tracts.

    \item \textbf{Figure S7: Multi-chromosome IBD overview.} Per-chromosome IBD segment comparison between hap-ibd (blue) and ibd-cli (coral) for all haplotype pairs in the validation subsets. Target regions shaded in yellow. (See Figure~\ref{fig:multi_chr_overview_supp}.)

    \item \textbf{Figure S8: All-chromosomes summary.} Aggregated view of hap-ibd vs.\ ibd-cli segments across chr10, chr11, chr12 in a three-row panel. Target regions shaded in green. (See Figure~\ref{fig:all_chr_summary_supp}.)

    \item \textbf{Figure S9: hap-ibd landscape.} Characterization of the hap-ibd ground truth across chr10, chr11, chr12, showing segment density, LOD distributions, and population composition. (See Figure~\ref{fig:hapibd_landscape_supp}.)

    \item \textbf{Figure S10: hap-ibd summary statistics.} Per-chromosome hap-ibd summary: segment count, mean length, LOD distribution, and population breakdown. (See Figure~\ref{fig:hapibd_stats_supp}.)

    \item \textbf{Figure S11: RFMix admixture summary.} Distribution of 3-way admixture proportions (AMR, EUR, AFR) across 43 AMR query samples on chr10/11/12. (See Figure~\ref{fig:rfmix_admixture_supp}.)

    \item \textbf{Figures S12a--S12i: Ancestry comparison paintings (chr10).} Side-by-side chromosome paintings comparing RFMix and ancestry-cli local ancestry calls for 9 AMR samples on chr10. Colors represent EUR (blue), AFR (red), and AMR (green) ancestry. Pangenome-based calls are shown only in the 1~Mb regions with available similarity data (5--6~Mb and 20--21~Mb). (See Figures~\ref{fig:ancestry_comparison_chr10_start}--\ref{fig:ancestry_comparison_chr10_end}.)

    \item \textbf{Figures S13a--S13i: Ancestry comparison paintings (multi-chromosome combined).} Combined chromosome paintings for all 9 AMR samples showing RFMix vs.\ ancestry-cli ancestry calls across chr10, chr11, and chr12 in a single view. Validated regions (6 total: chr10:5--6~Mb, chr10:20--21~Mb, chr11:25--26~Mb, chr11:70--71~Mb, chr12:5--6~Mb, chr12:50--51~Mb) are highlighted. Samples span a range of admixture profiles across the HPRC AMR cohort (HG00733, HG01071, HG01109, HG01175, HG01192, HG01928, HG01934, HG01943, NA19776). (See Figures~\ref{fig:ancestry_combined_HG00733}--\ref{fig:ancestry_combined_NA19776}.)

    \item \textbf{Figure S14: Robust vs.\ non-robust IBD statistics.} Four-panel comparison of IBD pair ranking under mean, trimmed mean (5\%), trimmed mean (10\%), and median statistics. The chr10 outlier pair (rank 103/120 by mean) resolves to rank 10/120 by trimmed mean and 6/120 by median. (See Figure~\ref{fig:robust_comparison}.)

    \item \textbf{Figure S15: Chr10 outlier diagnostic.} Three-panel diagnosis of the chr10 challenging case (HG01175\texttt{\#2} vs.\ HG01192\texttt{\#1}): per-window identity distribution showing 5 outlier windows, comparison with other IBD pairs, and rank improvement under robust statistics.

    \item \textbf{Figure S16: Window size sensitivity analysis.} Four-panel analysis of IBD detection robustness across effective window sizes (10--500~kb) under mean, trimmed mean (5\%), and median statistics. Detection remains stable at 10--100~kb; larger windows lose small regions due to insufficient data. (See Figure~\ref{fig:window_sensitivity}.)

    \item \textbf{Figure S17: IBD boundary precision.} Four-panel analysis of hap-ibd vs.\ pangenome boundary offsets across all 11 IBD pairs and 6 regions. Most hap-ibd segments extend beyond the pangenome analysis window; where boundaries are observable, end offsets are 20--31~kb.
\end{enumerate}

\begin{figure}[ht]
\centering
\includefigure[width=\textwidth,height=0.85\textheight,keepaspectratio]{../../scripts/validation/ibd/plots/chr10_ibd_overview.png}
\caption*{\textbf{Figure S7a}: Chr10 IBD overview. Horizontal bars show IBD segments per haplotype pair, comparing hap-ibd (blue) and ibd-cli (coral). Yellow shaded regions indicate the two 15~Mb target windows analyzed with pangenome IBS data.}
\end{figure}

\begin{figure}[ht]
\centering
\includefigure[width=\textwidth]{../../scripts/validation/ibd/plots/chr11_ibd_overview.png}
\caption*{\textbf{Figure S7b}: Chr11 IBD overview, formatted as S7a.}
\end{figure}

\begin{figure}[ht]
\centering
\includefigure[width=\textwidth]{../../scripts/validation/ibd/plots/chr12_ibd_overview.png}
\caption*{\textbf{Figure S7c}: Chr12 IBD overview. The single true positive detection (HG00738\#2 vs.\ HG01071\#1 at 117--125~Mb) is visible as a coral bar overlapping the hap-ibd blue bar in target region 2.}
\label{fig:multi_chr_overview_supp}
\end{figure}

\begin{figure}[ht]
\centering
\includefigure[width=\textwidth]{../../scripts/validation/ibd/plots/all_chromosomes_summary.png}
\caption{All-chromosomes IBD summary. Three-row panel showing hap-ibd vs.\ ibd-cli segments across chr10, chr11, and chr12. Green shaded areas indicate target regions. Total segment counts per method are annotated.}
\label{fig:all_chr_summary_supp}
\end{figure}

\begin{figure}[ht]
\centering
\includefigure[width=\textwidth]{../../scripts/validation/ibd/plots/hapibd_landscape.png}
\caption{hap-ibd ground truth landscape across chr10, chr11, chr12. Shows segment density along each chromosome, population composition of IBD sharing, and the chr10 hotspot at 45--55~Mb.}
\label{fig:hapibd_landscape_supp}
\end{figure}

\begin{figure}[ht]
\centering
\includefigure[width=\textwidth]{../../scripts/validation/ibd/plots/hapibd_summary_stats.png}
\caption{hap-ibd summary statistics. Per-chromosome segment counts, mean length, LOD score distribution, and population breakdown for the 223 HPRC sample pairs.}
\label{fig:hapibd_stats_supp}
\end{figure}

\begin{figure}[ht]
\centering
\includefigure[width=\textwidth]{../../scripts/validation/ibd/plots/chr10_identity_profiles.png}
\caption{Chr10 identity profiles for the top haplotype pairs in target regions. Pangenome IBS identity (y-axis, range 0.95--1.005) plotted along genomic position (x-axis) for individual pairs, overlaid with IBD segment calls from hap-ibd (blue) and ibd-cli (coral).}
\label{fig:chr10_identity_profiles}
\end{figure}

\begin{figure}[ht]
\centering
\includefigure[width=\textwidth]{../../scripts/validation/ibd/plots/chr11_identity_profiles.png}
\caption{Chr11 identity profiles, formatted as Figure~\ref{fig:chr10_identity_profiles}.}
\label{fig:chr11_identity_profiles}
\end{figure}

\begin{figure}[ht]
\centering
\includefigure[width=\textwidth]{../../scripts/validation/ibd/plots/chr12_identity_profiles.png}
\caption{Chr12 identity profiles. The validated IBD pair (HG00738\#2 vs.\ HG01071\#1) shows sustained identity at 0.9999--1.0 across the 117--125~Mb detection region (8~Mb, 800 windows).}
\label{fig:chr12_identity_profiles}
\end{figure}

\section{Validation Reference Panel}
\label{sec:supp_panel}

\begin{table}[ht]
\centering
\caption{Ancestry validation reference panel composition for the chr20 baseline (2-way: EUR/AFR). All 79 individuals are present in both the HPRC pangenome and the RFMix reference panel (100\% overlap). 16 AMR individuals serve as query samples.}
\label{tab:ref_panel}
\begin{tabular}{lcc}
\hline
\textbf{Population} & \textbf{Individuals} & \textbf{Haplotypes} \\
\hline
EUR & 30 & 60 \\
AFR & 49 & 98 \\
\hline
\textbf{Total reference} & 79 & 158 \\
AMR (query) & 16 & 32 \\
\hline
\end{tabular}
\end{table}

\begin{table}[ht]
\centering
\caption{Ancestry validation reference panel for chr10/11/12 (3-way: EUR/AFR/AMR). For RFMix, 1,486 reference haplotypes from the full phased VCFs were used. For ancestry-cli, the pangenome reference panel consists of 31 HPRC individuals whose assemblies are in the pangenome graph.}
\label{tab:ref_panel_3way}
\begin{tabular}{lcc}
\hline
\textbf{Population} & \textbf{Individuals} & \textbf{Haplotypes} \\
\hline
EUR & 8 & 16 \\
AFR & 8 & 16 \\
AMR (reference) & 15 & 30 \\
\hline
\textbf{Total reference} & 31 & 62 \\
AMR (query) & 9 & 18 \\
\hline
\end{tabular}
\end{table}

\section{Population Sample Sizes}
\label{sec:supp_samples}

\begin{table}[ht]
\centering
\caption{HPRC population sample sizes used in this study. Of 227 pangenome samples, 223 are present in the phased VCFs for chr10/11/12 (4 missing: HG01123, HG02486, HG02559, HG03471).}
\label{tab:sample_sizes}
\begin{tabular}{lccc}
\hline
\textbf{Population} & \textbf{Pangenome} & \textbf{In VCF} & \textbf{Haplotypes (VCF)} \\
\hline
AFR & 67 & 63 & 126 \\
EUR & 30 & 30 & 60 \\
EAS & 50 & 43 & 86 \\
CSA & 36 & 45 & 90 \\
AMR & 44 & 42 & 84 \\
\hline
\textbf{Total} & 227 & 223 & 446 \\
\hline
\end{tabular}
\end{table}

The 223 validated samples span all five continental populations, enabling population-stratified concordance analysis across intra- and inter-population pair categories (see Validation Design in the main text).

\section{Software Versions}
\label{sec:supp_software}

\begin{table}[ht]
\centering
\caption{Software versions used in this study.}
\label{tab:software}
\begin{tabular}{llp{8cm}}
\hline
\textbf{Tool} & \textbf{Version} & \textbf{Purpose} \\
\hline
impopk & 0.2.0 & IBD/IBS/ancestry/Jacquard analysis \\
impg & 0.2.0 & Pangenome similarity queries \\
AGC & 3.0 & Assembled genome storage and retrieval \\
minigraph-cactus & 2.0 & Pangenome graph construction \\
hap-ibd & 1.0 (15Jun23.92f) & VCF-based IBD detection (validation) \\
RFMix & 2.03 & Local ancestry inference (validation) \\
Rust & 1.70+ & Implementation language \\
rayon & 1.8 & Parallel processing \\
\hline
\end{tabular}
\end{table}

\section{3-Way Ancestry Validation Figures}
\label{sec:supp_ancestry_3way}

\begin{figure}[ht]
\centering
\includefigure{../../results/comparison_plots/rfmix_admixture_summary.png}
\caption{RFMix 3-way admixture proportions across 43 AMR query samples (chr10, chr11, chr12). Bar chart showing per-sample proportions of AMR (green), EUR (blue), and AFR (red) ancestry. Mean admixture: AMR = 71.1\%, EUR = 22.0\%, AFR = 6.8\%.}
\label{fig:rfmix_admixture_supp}
\end{figure}

\begin{figure}[ht]
\centering
\includefigure{../../results/comparison_plots/ancestry_HG00733_chr10.png}
\caption{Ancestry comparison: HG00733 chr10. Side-by-side RFMix (top) vs.\ ancestry-cli (bottom) chromosome painting. Ancestry-cli calls shown at 5--6~Mb and 20--21~Mb windows where similarity data is available. HG00733 is a Puerto Rican individual with predominantly AMR/EUR ancestry on chr10 (RFMix: AMR = 33.1\%, EUR = 66.9\%).}
\label{fig:ancestry_comparison_chr10_start}
\end{figure}

\begin{figure}[ht]
\centering
\includefigure{../../results/comparison_plots/ancestry_HG01109_chr10.png}
\caption{Ancestry comparison: HG01109 chr10. This sample has the highest AFR proportion in the validation set (genome-wide: AFR = 46.0\%, AMR = 39.5\%, EUR = 14.5\%), providing a strong test of minority ancestry detection.}
\end{figure}

\begin{figure}[ht]
\centering
\includefigure{../../results/comparison_plots/ancestry_HG01175_chr10.png}
\caption{Ancestry comparison: HG01175 chr10. An AMR individual providing additional validation of ancestry-cli performance across diverse admixture patterns in the HPRC panel.}
\label{fig:ancestry_comparison_chr10_end}
\end{figure}

% Multi-chromosome combined ancestry paintings (S13)
\begin{figure}[ht]
\centering
\includefigure{../../results/comparison_plots/ancestry_HG00733_combined.png}
\caption{Multi-chromosome ancestry comparison: HG00733 combined (chr10, chr11, chr12). RFMix (top) vs.\ ancestry-cli (bottom) showing all validated regions across three chromosomes.}
\label{fig:ancestry_combined_HG00733}
\end{figure}

\begin{figure}[ht]
\centering
\includefigure{../../results/comparison_plots/ancestry_HG01109_combined.png}
\caption{Multi-chromosome ancestry comparison: HG01109 combined. This sample has the highest AFR proportion in the validation set (genome-wide: AFR = 46.0\%), providing a cross-chromosome test of minority ancestry detection.}
\label{fig:ancestry_combined_HG01109}
\end{figure}

\begin{figure}[ht]
\centering
\includefigure{../../results/comparison_plots/ancestry_HG01175_combined.png}
\caption{Multi-chromosome ancestry comparison: HG01175 combined (chr10, chr11, chr12). An AMR individual showing ancestry-cli calls compared to RFMix across all validated regions.}
\label{fig:ancestry_combined_HG01175}
\end{figure}

\begin{figure}[ht]
\centering
\includefigure{../../results/comparison_plots/ancestry_HG01192_combined.png}
\caption{Multi-chromosome ancestry comparison: HG01192 combined (chr10, chr11, chr12). An AMR individual showing ancestry-cli calls compared to RFMix across all validated regions.}
\label{fig:ancestry_combined_HG01192}
\end{figure}

\begin{figure}[ht]
\centering
\includefigure{../../results/comparison_plots/ancestry_HG01928_combined.png}
\caption{Multi-chromosome ancestry comparison: HG01928 combined (chr10, chr11, chr12). An AMR individual showing ancestry-cli calls compared to RFMix across all validated regions.}
\label{fig:ancestry_combined_HG01928}
\end{figure}

\begin{figure}[ht]
\centering
\includefigure{../../results/comparison_plots/ancestry_HG01071_combined.png}
\caption{Multi-chromosome ancestry comparison: HG01071 combined. Predominantly AMR individual (genome-wide: AMR = 78.9\%, EUR = 20.2\%, AFR = 0.9\%).}
\label{fig:ancestry_combined_HG01071}
\end{figure}

\begin{figure}[ht]
\centering
\includefigure{../../results/comparison_plots/ancestry_HG01934_combined.png}
\caption{Multi-chromosome ancestry comparison: HG01934 combined (chr10, chr11, chr12). An AMR individual showing ancestry-cli calls compared to RFMix across all validated regions.}
\label{fig:ancestry_combined_HG01934}
\end{figure}

\begin{figure}[ht]
\centering
\includefigure{../../results/comparison_plots/ancestry_HG01943_combined.png}
\caption{Multi-chromosome ancestry comparison: HG01943 combined (chr10, chr11, chr12). An AMR individual showing ancestry-cli calls compared to RFMix across all validated regions.}
\label{fig:ancestry_combined_HG01943}
\end{figure}

\begin{figure}[ht]
\centering
\includefigure{../../results/comparison_plots/ancestry_NA19776_combined.png}
\caption{Multi-chromosome ancestry comparison: NA19776 combined. Three-way admixed individual with substantial AFR component (genome-wide: AMR = 46.9\%, EUR = 30.5\%, AFR = 22.5\%).}
\label{fig:ancestry_combined_NA19776}
\end{figure}

% IBD robust statistics and sensitivity figures (S14--S17)
\begin{figure}[ht]
\centering
\includefigure{../../scripts/validation/ibd/plots/robust_comparison.png}
\caption{Robust vs.\ non-robust IBD statistics (S14). Four-panel comparison of IBD pair ranking under mean, trimmed mean (5\%), trimmed mean (10\%), and median statistics. The chr10 outlier pair (rank 103/120 by mean) resolves to rank 10/120 by trimmed mean and 6/120 by median, demonstrating that robust statistics recover correct ranking despite local identity drops.}
\label{fig:robust_comparison}
\end{figure}

\begin{figure}[ht]
\centering
\includefigure[width=\textwidth]{../../scripts/validation/ibd/plots/chr10_outlier_diagnostic.png}
\caption{Chr10 outlier diagnostic (S15). Three-panel diagnosis of HG01175\texttt{\#2} vs.\ HG01192\texttt{\#1}: (left) per-window identity distribution with 5 outlier windows at identity $<$ 0.9, (center) comparison with other IBD pairs, (right) rank improvement under robust statistics (103/120 by mean $\rightarrow$ 6/120 by median).}
\label{fig:chr10_outlier_diagnostic}
\end{figure}

\begin{figure}[ht]
\centering
\includefigure{../../scripts/validation/ibd/plots/window_sensitivity.png}
\caption{Window size sensitivity analysis (S16). Four-panel analysis of IBD detection robustness across effective window sizes (10--500~kb) under mean, trimmed mean (5\%), and median statistics. Detection remains stable at 10--100~kb; larger windows lose resolution for small regions due to insufficient data points.}
\label{fig:window_sensitivity}
\end{figure}

\begin{figure}[ht]
\centering
\includefigure[width=\textwidth]{../../scripts/validation/ibd/plots/boundary_precision.png}
\caption{IBD boundary precision (S17). Comparison of hap-ibd segment boundaries vs.\ pangenome analysis window boundaries across 11 IBD pairs. Most hap-ibd segments extend beyond the analysis window (start/end at edge); where observable, end boundary offsets are 20--31~kb.}
\label{fig:boundary_precision}
\end{figure}

\section{Extended Discussion}
\label{sec:supp_extended_discussion}

The main-text Discussion highlights key findings; this section provides extended analysis on specific topics.

\subsection{Population-Specific Modeling}
\label{sec:supp_pop_modeling}

A key technical contribution is the explicit modeling of population-specific nucleotide diversity in the HMM emission distributions. The difference in $\pi$ between populations (AFR: 0.00125, EUR: 0.00085) directly affects the separation between IBD and non-IBD distributions. Without population-specific calibration, methods tuned to European diversity levels will be overly sensitive in African populations (detecting spurious IBD) or insufficiently sensitive when applied to populations with lower diversity.

Our population-adaptive transition model further accounts for the different IBD landscapes across populations. African populations have shorter and rarer IBD segments due to their larger effective population size and deeper coalescence times. By scaling transition probabilities accordingly (0.7$\times$ expected segment length and 0.3$\times$ entry probability for AFR), the model maintains appropriate sensitivity and specificity across diverse populations.

The MAP-regularized EM and BIC model selection provide additional robustness for challenging cases. In African populations, the IBD and non-IBD distributions overlap heavily, causing standard k-means clustering to fail. The MAP regularization anchors the emission estimates near population priors, preventing the component means from collapsing to a single mode. BIC model selection then guards against over-fitting by rejecting the two-component model when the data does not support it.

A notable consequence of operating on assembly identity rather than genotyped markers is the inversion of the population performance hierarchy. SNP array-based methods achieve their highest power in European populations (where arrays were designed) and lowest in African populations. In contrast, IBD detection SNR from pangenome identity is highest in African populations (1.06) and lowest in East Asian populations (0.69), because the IBD--non-IBD identity gap scales with nucleotide diversity $\pi$, which is greatest in populations with the deepest coalescent times.

\subsection{Detailed Error Decomposition}
\label{sec:supp_error_decomposition}

Diagnostic analysis of simulated ancestry errors reveals three distinct sources: (i) \emph{weak signal} (54--72\% of errors)---windows where the identity difference between correct and incorrect populations is $<$0.0005; (ii) \emph{wrong signal} (37--40\%)---windows where the closest reference haplotype belongs to the wrong population due to incomplete lineage sorting or shared deep ancestry; and (iii) \emph{HMM miss} (5--12\%)---windows where signal was present but smoothed away by the transition model. The 2.05\% residual error decomposes into five sources: short tracts below HMM detection ($\sim$0.7\,pp), panel coverage gaps ($\sim$0.5\,pp), emission bias from within-population substructure ($\sim$0.3\,pp), boundary precision limited by Viterbi inertia ($\sim$0.3\,pp), and an irreducible floor from ILS and convergent mutation ($\sim$0.25\,pp). The empirical ceiling is approximately 99.75\%.

\subsection{Prediction Epistemology}
\label{sec:supp_prediction}

Systematic evaluation of analytical predictions reveals a sharp epistemological boundary. Descriptive predictions---characterizing equilibrium properties (Fisher information content, spectral equivalence, transition dominance)---achieve 100\% accuracy (3/3 verified). Prescriptive predictions---forecasting intervention effects---achieve 0\% accuracy (0/8 verified; Fisher exact $p = 0.006$). The root cause is auto-configuration: component modifications cascade through parameter re-estimation, producing indirect effects that dominate predicted direct effects. This pattern parallels the Lucas critique in economics.

\subsection{Confidence-Concordance Tradeoff}
\label{sec:supp_confidence}

All pairwise emission variants increase model confidence but decrease concordance with RFMix on HPRC data: hybrid pairwise ($+$3.3\,pp confidence, $-$5.20\,pp concordance), adaptive pairwise ($+$1.3\,pp, $-$4.89\,pp), and agreement-based pairwise ($+$0.5\,pp, $-$3.05\,pp). The root cause is EUR--AMR genealogical entanglement: most AMR haplotype windows are genealogically EUR or AFR, so pairwise emissions detect genealogical truth rather than the population-of-origin label assigned by RFMix. The tradeoff is task-dependent: standard emissions are optimal for LAI concordance with admixed reference panels, while pairwise emissions are preferred for eGRM estimation.

\subsection{Segmental Duplication Masking}
\label{sec:supp_sd_masking}

Chromosome 1 contains $\sim$6\% segmental duplications ($\sim$525 regions, $\sim$16\,Mbp) compared to $\sim$2\% for chr12. SD masking was hypothesized to explain the 11.7\,pp accuracy gap between chr12 (76.45\%) and chr1 (68.12\%). However, empirical validation definitively disproved this: masking SD regions \emph{hurts} accuracy by 1--7\,pp across both backends. SD windows carry more ancestry signal than average, likely because paralogous mappings increase effective reference panel size. The identity floor (\texttt{--identity-floor 0.9}) already subsumes SD artifact removal.

\subsection{Automatic Parameter Configuration Details}
\label{sec:supp_autoconfig}

Optimal ancestry parameters differ substantially between datasets---applying simulation-tuned parameters to HPRC data incurs a 9.8\,pp penalty. The auto-configure system eliminates this by computing $D_{\min}$ (minimum Cohen's $d$ across population pairs) from a calibration scan and deriving parameters from this single statistic. An attempt to improve robustness via Tukey fence outlier filtering proved counterproductive: HPRC windows have near-zero IQR, causing the threshold to capture 11--13\% of normal windows and inflating $D_{\min}$ by 10--20$\times$. Cross-chromosome parameter transfer improves accuracy by 4.02\,pp (MSE $\approx$ Bias$^2 \gg$ Var), and works via the \texttt{--save-params}/\texttt{--load-params} interface.

\subsection{Logit-Space IBD Emissions: Theory and Negative Result}
\label{sec:supp_logit}

Mapping identity through the logit transform improves Gaussian model adequacy (skewness 4.1$\times$ reduced, boundary leakage 11.7\%$\to$0\%), with a net $\sim$5\% per-window efficiency gain. However, empirical evaluation on HPRC data shows a regression ($-$0.226 F1): the logit amplifies the K=0 spike ($\sim$22\% of IBD windows) into a discrete mass at the cap, creating bimodal emissions that violate Gaussian assumptions worse than the original distribution.

\subsection{Asymmetric Detection: the \texorpdfstring{$\pi^{3/2}$}{pi\^{}(3/2)} Law}
\label{sec:supp_pi_law}

Minority ancestry detection degrades as $\pi_j^{3/2}$ due to three compounding penalties: fewer tracts ($\propto \pi_j$), higher prior penalty ($\propto 1/\sqrt{\pi_j}$), and noisier emissions ($\propto 1/\sqrt{\pi_j}$). At 50/30/20 admixture, AMR at 20\% has $\sim$1/4 the detection power of EUR at 50\%. A ``help-where-least-needed'' paradox follows: emission improvements disproportionately benefit the majority population. Optimal panel allocation provides $<$0.3\,pp gain at current panel sizes.

\subsection{Spectral Gap and Minimum Detectable Tracts}
\label{sec:supp_spectral}

The ancestry HMM spectral gap $\gamma = sK/(K-1)$ governs the correlation length $\xi \approx 1{,}700$ windows ($\sim$17\,Mb) and mixing time $\sim$100\,Mb. Population-specific minimum detectable tracts: AFR $\sim$360\,kb, EUR/AMR $\sim$2.3\,Mb. This $\sim$6$\times$ difference reflects the $F_{ST}$ hierarchy and explains per-ancestry detection patterns.

\subsection{Instructive Negative Results}
\label{sec:supp_negative}

Posterior decoding produces results identical to Viterbi because posteriors rarely favor minority ancestry. Emission normalization (z-score) is catastrophic ($-$83\,pp), because the systematic EUR--AFR similarity difference \emph{is} the discriminative signal. A multi-feature oracle testing all combinations of seven alignment features achieved accuracy exactly equal to the majority-class baseline with zero minority recall, proving that per-window information content is fundamentally insufficient---the HMM's concordance is achieved entirely through temporal smoothing.

\section{Extended Results}
\label{sec:supp_extended_results}

This section contains detailed analyses, parameter sweeps, and per-pair breakdowns that support the main-text results.

\subsection{hap-ibd Ground Truth Landscape}
\label{sec:supp_hapibd_ground_truth}

hap-ibd was run on chr10, chr11, and chr12 using phased VCFs (4,091 samples each), filtered to 223 HPRC sample pairs. Chr10 is dominated by a massive hotspot at 45--55\,Mb (5,222 of 5,491 segments, 95\%), analogous to the chr15q11--q13 hotspot. Outside this hotspot, chr10 shows only 269 segments---comparable to chr11 (264) and chr12 (348). Total segments across three chromosomes: 6,103, with 505 high-confidence segments (LOD~$\geq$~3).

\subsection{IBD Parameter Sweep (chr20)}
\label{sec:supp_ibd_params}

Nine parameter configurations were evaluated on chr20 (44 sample pairs). The best configuration uses identity floor~=~0.9 with 2\,Mb minimum segment length and 10 Baum-Welch iterations (F1~=~28.3\%). Lower identity floors allow alignment-gap noise; higher floors remove informative windows. Baum-Welch training provides $+$110\% F1 improvement vs.\ no training. The 2\,Mb minimum length balances precision (12 segments) against recall (3\,Mb loses recall, 1\,Mb produces 140--220 segments).

\subsection{Pangenome Identity Distribution Challenge}
\label{sec:supp_identity_dist}

The moderate IBD concordance reflects a fundamental difference in identity distributions. For a representative pair (HG00280 vs.\ HG00323), 62\% of non-IBD windows have pangenome identity $>$0.99, making them indistinguishable from true IBD windows. The distribution is strongly bimodal (identity near 0 or near 1), with standard deviation ($\sim$0.29) dwarfing the IBD signal ($\Delta\mu \approx 0.003$) by $\sim$100$\times$. While the raw IBD/non-IBD separation is actually larger in pangenome data ($\sim$0.003 vs.\ $\sim$0.00055 for VCF), the dramatically higher variance renders the signal-to-noise ratio approximately 55$\times$ worse.

\subsection{Performance Ceiling: Eight Approaches Converge}
\label{sec:supp_ceiling}

Eight algorithmic improvements beyond the baseline were evaluated: distance-aware transitions, genetic-map transitions, logit-transform emissions, background filtering, z-score normalization, population-adaptive transitions, coverage-ratio auxiliary emission, and region exclusion. None improved upon the baseline F1~=~28.3\%. Distance-aware and genetic-map transitions produce identical calls (emission bottleneck, not transition). Logit-transform is counterproductive (identity~=~1.0 outliers). Z-score normalization is incompatible with HMM priors. Region exclusion removes true positives alongside artifacts (F1 drops to 15.9\%).

\subsection{False Positive Hotspots and Missed Segments}
\label{sec:supp_fp_analysis}

Two false positive hotspot regions dominate: chr20:45--50\,Mb (41\% of FPs) and chr20:35--40\,Mb (16\%), together accounting for 57\% of all false positives within 23\% of the chromosome. The only reliable discriminator between true and false positive segments is segment length (TP mean 3.43\,Mb vs.\ FP mean 2.35\,Mb), with 82\% of FPs in the 2.0--2.5\,Mb bin.

\subsection{Scaled IBD Validation (2,850 Pairs)}
\label{sec:supp_scaled_ibd_summary}

At full scale on chr20 (76 HPRC samples, 2,850 pairs), our pipeline detects 293 IBD segments across 258 pairs; hap-ibd reports 17 segments across 16 pairs. Pair-level: 11 TP, 247 FP, 5 FN, precision 4.3\%, recall 68.8\%. EUR pairs have the best recall (88.9\%, 8/9 detected). AMR--EUR cross-population pairs are FP-dominant (all 90 detections are false positives). The best operating points are 3.0\,Mb minimum length (segment F1~=~43.2\%) or LOD~$\geq$~300 (F1~=~44.4\%).

\subsection{Haplotype-Level Concordance and Boundaries}
\label{sec:supp_haplotype}

Only 4/11 TP pairs have matching haplotype labels between HPRC assemblies and VCF phasing, confirming that assembly haplotype labels are arbitrary. Sample-level concordance is the appropriate cross-tool metric. Segment boundary accuracy is 3.35\,Mb (start) and 1.78\,Mb (end) on average, reflecting gradual identity transitions at boundaries.

\subsection{Structural Variant Enrichment}
\label{sec:supp_sv}

Pangenome IBD detection performs \emph{worse} in structurally variable regions: Q4 (highest SV) has zero true positive detections with mean identity 0.864---far below the IBD range ($>$0.999). Q1 (lowest SV) achieves the best F1 (0.320). The hypothesis that pangenome identity would capture IBD in SV regions is not supported: structural variation degrades alignment quality.

\subsection{2-Way Ancestry Detailed Analysis (chr20)}
\label{sec:supp_2way_ancestry}

With the \texttt{max} emission model on chr20 (HG00733, EUR/AFR), overall concordance reaches 93.2\%. EUR F1~=~0.965; AFR F1~=~0.158. SNR analysis shows that in true AFR windows, EUR references have \emph{higher} similarity in 75.2\% of cases. A multi-feature oracle (seven alignment features, all 2--3 feature combinations) achieves accuracy exactly equal to the majority-class baseline with zero AFR recall, proving per-window features are non-informative. Multi-sample validation across 15 AMR individuals yields mean concordance 93.4\% (range 90.2--95.5\%). Seven approaches for improving minority detection were tested; only increasing $p_{\text{switch}}$ to 0.02 consistently helps (AFR F1: 0.158 $\rightarrow$ 0.28, at $-$3.1\,pp overall cost).

\subsection{3-Way Ancestry Detailed Analysis}
\label{sec:supp_3way_detail}

Emission model sweep across 13 configurations shows top-3 averaging (76.2\%) as optimal. Top-2 produces identical results; top-10 degrades to 60.9\% (signal dilution). Genetic-map transitions, smoothing, and posterior decoding provide no additional improvement. Per-region concordance ranges from 66.7\% (chr10:5--6\,Mb) to 88.9\% (chr11:25--26\,Mb). Confusion matrix: 900 AMR$\rightarrow$AFR and 1,800 AMR$\rightarrow$EUR misclassifications dominate.

\subsection{Multi-Chromosome IBD Detailed Analysis}
\label{sec:supp_multi_chr_ibd}

Ranking-based evaluation across six regions on chr10/11/12: 10/11 IBD pairs rank in top 10\% by mean identity; 6 rank \#1. The single outlier (HG01175 vs.\ HG01192, chr10) has 5 artifact windows dragging the mean; robust statistics (trimmed mean, median) achieve 11/11 (100\%) detection. Detection is stable across 10--100\,kb effective window sizes. At region-specific thresholds, 5/6 regions achieve perfect precision--recall. Segment length discriminates TP ($>$3.8\,Mb) from FP ($<$2.5\,Mb). Identity gap $\Delta\mu = 0.003 \pm 0.0004$ is reproducible across all chromosomes.

\section{Extended Methods}
\label{sec:supp_methods}

\subsection{IBS Computation Details}
\label{sec:supp_ibs_details}

\paragraph{Sliding window formulation.}
Given a target genomic region $[s, e]$ and window size $W$, the region is partitioned into non-overlapping windows $\mathcal{W} = \{[s + kW, \min(s + (k+1)W - 1, e)] \mid k = 0, 1, \ldots\}$. Windows are processed in parallel using the \texttt{rayon} work-stealing thread pool.

\paragraph{Sequence identity via impg.}
For each window, \texttt{impg similarity} is invoked on the corresponding reference coordinates using the pangenome reference coordinate system (\texttt{ref\#0\#chrom:start-end}), projecting alignments through the pangenome graph structure. The tool queries the alignment (PAF or 1aln format) against an AGC sequence store.

\paragraph{Jacquard coefficient estimation.}
For two diploid individuals $A$ and $B$ with haplotypes $(a_1, a_2)$ and $(b_1, b_2)$, the IBS sharing pattern across all four haplotype pairs $\{(a_i, b_j)\}$ at each locus determines the condensed identity state ($\Delta_1$ through $\Delta_9$), enabling estimation of Jacquard delta coefficients without phased genotype calls.

\paragraph{Batch processing.}
For whole-genome analysis, multiple genomic regions can be specified via a BED file (0-based, half-open coordinates converted to 1-based inclusive internally). Each region is processed sequentially, with windows parallelized across available threads.

\subsection{IBD HMM Details}
\label{sec:supp_ibd_details}

\paragraph{Non-IBD emission variance.}
The variance accounts for Poisson sampling of differences within a window, with a correction for linkage disequilibrium:
$\sigma_0 = \sqrt{\pi \cdot c_{\text{LD}} / W}$
where $W$ is the window size and $c_{\text{LD}} \approx 3$ is the LD correction factor. The Cram\'er-Rao lower bound for boundary precision is $\geq 1.35W \approx 13.5$\,kb for 10\,kb windows, close to the observed 15--32\,kb precision.

\paragraph{Population-adaptive transition scaling.}
AFR: $L' = 0.7L$, $p'_{\text{enter}} = 0.3\,p_{\text{enter}}$ (fewer, shorter IBD segments due to deeper coalescence). EAS: $L' = 1.1L$ (slightly longer segments). Inter-population: $L' = 0.5L$, $p'_{\text{enter}} = 0.1\,p_{\text{enter}}$ (cross-population IBD is rare).

\paragraph{Distance-dependent transitions.}
When consecutive windows are non-uniformly spaced, transitions are modeled using a continuous-time Markov chain. The per-base-pair rate is $\lambda = -\ln(1 - p_{\text{switch}})/W$, and the transition probability over distance $d$ is $P(\text{switch} \mid d) = 1 - e^{-\lambda d}$, clamped to $[10^{-10}, 1 - 10^{-10}]$.

\paragraph{Recombination-rate-aware transitions.}
Given a genetic map, genetic distance $d_{\text{cM}}$ between window midpoints is computed by linear interpolation. The Haldane map function gives recombination probability $r = \frac{1}{2}(1 - e^{-2d_M})$ where $d_M = d_{\text{cM}}/100$. Transition rates are scaled by $s = d_{\text{cM}} / d_{\text{cM}}^{\text{nominal}}$:
$P(\text{enter IBD} \mid d) = 1 - e^{-\lambda_{\text{enter}} \cdot s}$,
$P(\text{exit IBD} \mid d) = 1 - e^{-\lambda_{\text{exit}} \cdot s \cdot (1 + r)}$.

\paragraph{Initial emission estimation.}
K-means clustering (up to 30 iterations) assigns observations to two groups. If separation is insufficient ($\Delta\mu \leq 0.0005$), a two-component Gaussian mixture is fit via EM with MAP regularization (prior strength $\kappa$: AFR = 15, inter-pop = 10, others = 5). BIC model selection prevents overfitting.

\paragraph{Baum-Welch convergence analysis.}
The asymptotic convergence rate is $r \approx 4\rho(1-\rho)/(1-2\rho)^2$ where $\rho = \Phi(-\Delta\mu/2\sigma)$. For the IBD model (EUR: $\Delta\mu = 0.00055$, $\sigma \approx 0.00074$), $\rho \approx 0.36$ yields $r \approx 0.9$---each iteration reduces error by only $\sim$10\%. We recommend 20--30 iterations for IBD; the ancestry HMM converges faster ($r \approx 0.5$--$0.7$) because only transitions are re-estimated.

\paragraph{Biological bounds.}
After each Baum-Welch M-step: $\mu_0 \in [\mu_0^{\text{prior}} - 0.005, 0.9993]$, $\mu_1 \in [0.9990, 1.0]$, $\sigma_0 \in [0.0003, 0.005]$, $\sigma_1 \in [0.0002, 0.002]$, $\mu_0 < \mu_1$ (swap or reset if violated), $P(\text{enter}) \in [10^{-8}, 0.1]$, $P(\text{exit}) \in [0.001, 0.5]$.

\paragraph{Distance-aware Baum-Welch.}
When distance-aware mode is enabled, both the forward and backward passes use per-step distance-dependent transition matrices: $\xi_t(i,j) = \alpha_t(i) \cdot P_d(j|i) \cdot P(o_{t+1}|j) \cdot \beta_{t+1}(j) / P(\mathbf{O})$. The M-step re-estimates base rates per nominal window step, which are then scaled by distance during inference.

\paragraph{Quality score components.}
$Q \in [0, 100]$: posterior strength (0--40, mean $\gamma_t(1)$), posterior consistency (0--20, min/mean ratio), LOD evidence (0--30, LOD per window), segment length (0--10). Categories: $Q \geq 80$ (high), $Q \geq 50$ (medium), $Q < 20$ (likely FP).

\paragraph{Identity floor analysis.}
Pangenome-derived identity exhibits a bimodal distribution: windows either have identity near zero (alignment gaps) or near one (aligned regions). This bimodality introduces variance $\sigma \approx 0.29$ that overwhelms the IBD signal ($\Delta\mu \approx 0.003$). A floor of 0.98 removes gap artifacts while preserving genuine variation.

\subsection{Viterbi and Forward-Backward Details}
\label{sec:supp_inference_details}

The Viterbi algorithm operates in log-space: $v_t(s) = \max_{s'}[v_{t-1}(s') + \ln a_{s's}] + \ln P(o_t \mid s)$. Forward-backward posteriors:
\begin{align}
\alpha_t(s) &= \left[\textstyle\sum_{s'} \alpha_{t-1}(s') \cdot a_{s's}\right] \cdot P(o_t \mid s) \\
\beta_t(s) &= \textstyle\sum_{s'} a_{ss'} \cdot P(o_{t+1} \mid s') \cdot \beta_{t+1}(s') \\
\gamma_t(s) &= \alpha_t(s) \cdot \beta_t(s) / \textstyle\sum_{s'} \alpha_t(s') \cdot \beta_t(s')
\end{align}
All computations use log-sum-exp for numerical stability. Segment extraction uses run-length encoding with boundary refinement (extend at $\gamma > 0.5$, trim at $\gamma < 0.2$), minimum identity threshold (0.9995), maximum 1-window gap bridging, and minimum 3-window segment length.

\subsection{Ancestry HMM Details}
\label{sec:supp_ancestry_details}

\paragraph{Emission normalization.}
Optional per-population z-score normalization before softmax: $z_k(o_t) = (f_k(o_t) - \hat{\mu}_k)/\hat{\sigma}_k$, where $\hat{\mu}_k$ and $\hat{\sigma}_k$ are computed across all windows. Temperature then operates in z-score space (range 0.5--5.0).

\paragraph{Baum-Welch for ancestry.}
The switch probability is refined via EM over all sample sequences: E-step computes forward-backward probabilities and expected transition counts $\xi_t(i,j)$; M-step re-estimates the transition matrix, extracts the average off-diagonal mass as the updated switch probability, symmetrizes, and clamps to $[0.0001, 0.1]$.

\paragraph{Cross-validation.}
Model performance is evaluated via stratified $k$-fold cross-validation ($k = 5$): partition reference haplotypes by population, hold out one fold, train on remaining, and compute per-population precision, recall, F1.

\paragraph{Ancestry LOD scores.}
Per-window LOD: $\text{LOD}_t = \log_{10}[P(o_t \mid k^*) / P(o_t \mid k^{(2)})]$, where $k^*$ is the assigned and $k^{(2)}$ the second-best population. Segment LOD is the sum across segment windows.

\paragraph{Admixture proportion estimation.}
$\hat{q}_k = \sum_{i:a_i=k}(e_i - s_i) / \sum_i(e_i - s_i)$, where the sums are over ancestry segments. Additional statistics: ancestry switch count, mean and per-population tract lengths.

\subsection{Validation Parameter Details}
\label{sec:supp_validation_details}

\paragraph{IBD concordance metrics.}
Segment Jaccard: $J = |S_{\text{ours}} \cap S_{\text{theirs}}| / |S_{\text{ours}} \cup S_{\text{theirs}}|$. Precision and recall treat hap-ibd as ground truth. F1 is the harmonic mean. Per-window concordance classifies each 5\,kb window as IBD if $>$50\% covered. Matched segment length correlation uses reciprocal 50\% overlap criterion for segment matching.

\paragraph{hap-ibd parameters.}
hap-ibd v1.0 (15Jun23.92f), default parameters (min-seed = 2\,cM, min-output = 2\,cM) on phased VCFs (4,091 samples each). Results: chr10: 5,491 segments (hotspot at 45--55\,Mb contributes 95\%); chr11: 264 segments; chr12: 348 segments.

\paragraph{ibd-cli parameters.}
10\,kb windows via \texttt{impg}, six focused regions across chr10/11/12 selected to overlap high-LOD hap-ibd segments. Validation subsets of 8--16 samples selected by greedy optimization of hap-ibd pair coverage.

\paragraph{Ancestry concordance metrics.}
Per-population precision/recall/F1, switch point accuracy (fraction detected within $\pm k$ windows, mean distance), $N \times N$ confusion matrix, segment-level Jaccard.

\paragraph{RFMix parameters.}
RFMix v2.03, default parameters, 4 threads, $\sim$52 min per chromosome. Chr20 baseline: 2-way (EUR/AFR), 79 reference samples. Chr10/11/12: 3-way (EUR/AFR/AMR), 1,486 reference haplotypes.

\paragraph{ancestry-cli parameters.}
\texttt{max} emission model, \texttt{--estimate-params}, identity floor 0.9. Chr20: $\tau = 0.03$, $p_{\text{switch}} = 0.01$. Chr10/11/12: auto-estimated parameters, 240-sequence subset via \texttt{--subset-sequence-list}.

\paragraph{Statistical considerations.}
When one ancestry class dominates ($>$90\%), overall concordance is uninformative; per-population F1 and confusion matrices are more diagnostic. Disagreements between tools do not necessarily indicate errors, as the two use fundamentally different inputs. Boundary discordance of $\pm$1--2 windows is expected from finite window/CRF sizes.

\bibliographystyle{unsrtnat}
\bibliography{references}

\end{document}
